\scanpage{331}\mspage{247}
% l. 1–3
d'où un foncteur composé
\[
 C_j/\unsure{\varphi|\phi}(j)\unsure{\mathbin{-}|\longrightarrow}
 C_{j'}/\alpha^{\unsure{*|\times}}(\varphi(j))\unsure{\mathbin{-}|\longrightarrow}
 C_{j'}/\varphi(j'),
\]
avec transitivité \unsure{à|:} \unsure{Iso.|iso.} can. près etc.
% l. 4–8
On note que l'on a un isom. can.
\[
 \underline{\Gamma}(\mathscr D_{\varphi}/\deleted{\unsure{\mathscr D|J}}\unsure{J|\mathbb J})\simeq
 \underline{\Gamma}(\mathscr D/J)_{/\varphi}
\]
de sorte que le foncteur \inserted{fidèle} canonique (4.10.2.2.) s'identifie au foncteur
\[
 \underline{\Gamma}(\mathscr D_{\varphi}/J)\longrightarrow
 \prod_{j\in\operatorname{Ob}J}(\mathscr D_{\varphi})_j\,.\unsure{.} \tag*{(4.10.6)}
\]

% l. 9–11
Le corollaire 4.10.2 sera alors une conséquence du résultat plus général \deleted{\unsure{s}} suivant :

% l. 12–15
\uline{Corollaire} 4.10.3. Soit $\unsure{\mathcal E|\mathscr E}\longrightarrow J$ un\unsure{e} \deleted{catégorie} \inserted{foncteur} \unsure{cofibrant|cofibrante}, et $\unsure{\pi|\kappa}$ un cardinal infini. On suppose \deleted{$\mathcal E,J$ \unsure{équiv} \inserted{que $\mathcal E,J$ équiv. à de} petites catégories, et de plus.}
% l. 16–20

a) $\operatorname{Card}J\leq\pi$.

b) Les catégories fibres $\mathcal E_j$ sont équivalentes à des petites catégories, et grdes devant $\pi$.

c) \deleted{\unsure{Elles}} \unsure{Il existe} un ordinal \deleted{infini} \inserted{limite} $\alpha$, tel que $1^\circ)$ \deleted{$\beta<\alpha$ \unsure{induit que|implique}} $\operatorname{Card}\deleted{\beta}\alpha\leq\pi$, et
% archive sheet number: 330
